\section{Relativized reducts and admissible separation}
\label{sec:p03-relativized-admissible-separation}

Barwise's projective classes allow the reduct to have a smaller universe
than its witnessing expansion, not just fewer relation symbols
\cite[Part~II, Definition~1.1, pp.~248--249]{Barwise1974Axioms}.
Accordingly the ordinary interpolation theorem requires a further
argument before it gives separation for these classes. We retain the
countable admissible set \(A\), with \(\omega\in A\), of
\cref{def:p03-adm-setting}.

\begin{hbdefinition}{def:p03-adm-relativized-pc}{Relativized projective classes}
Let \(\sigma\in A\) be a relational signature. A class \(K\) of
\(\sigma\)-structures is projective in \(L_A\) if some expanded
signature \(\tau\in A\), distinguished unary predicate \(U\notin\sigma\),
and sentence \(\phi\in L_A(\tau)\) have the following property:
\(M\in K\) if and only if there is a model \(N\models\phi\)
whose induced \(\sigma\)-structure on \(U^N\)
is isomorphic to \(M\). Additional symbols and additional objects are
both allowed in \(N\). In this section empty universes are allowed,
as in Barwise's source convention; positive-arity functions on them
are empty functions, and a signature with a constant has no empty
structure. Nullary relations retain their truth values. A definable
domain can be named by \(U\) using
its defining equivalence. A class is \(\Delta^1_1\) in this sense if
both it and its complement among \(\sigma\)-structures are projective.
The Souslin--Kleene property says that each such complementary pair is
defined by a sentence of the logic itself.
\end{hbdefinition}

\begin{hblemma}{thm:p03-adm-exterior}{proved=annex:p3.adm.10}{Removing a pure exterior}
Let \(\theta\in L_A(\sigma\cup\{U\})\) be a sentence, with \(\sigma\)
relational and without constants. There is a sentence
\(\theta^\circ\in L_A(\sigma)\) such that the following holds.
For any \(\sigma\)-structure \(M\), including an empty one, adjoin a disjoint infinite set
of objects, interpret \(U\) as the original universe, and make every
\(\sigma\)-relation false on tuples containing an adjoined object.
Call the result \(M^+\). Then
\[
 M^+\models\theta\quad\Longleftrightarrow\quad M\models\theta^\circ.
\]
The same sentence works for every infinite cardinality of the adjoined
set. Its relation symbols already occur in \(\theta\), other than \(U\).
\end{hblemma}

\begin{hbtheorem}{thm:p03-adm-relativized-separation}{proved=annex:p3.adm.11}{Relativized separation}
Two disjoint classes of \(\sigma\)-structures projective in \(L_A\)
are separated by a sentence in \(L_A(\sigma)\): it holds throughout the
first class and fails throughout the second. In particular \(L_A\)
satisfies the Souslin--Kleene property with relativized reducts.
For finite \(\sigma\), the separator uses only finitely many symbols.
Function and constant signatures have the same conclusion after
relational graph translation and its inverse.
\end{hbtheorem}

This result concerns projective classes with possible auxiliary universes.
The Beth theorem in \cref{thm:p03-adm-beth} concerns uniqueness of a
relation on a fixed reduct. Neither is replaced by an unqualified
identification of those two notions of definability.
